\chapter{Strategie di Esplorazione}

Le cinque strategie implementano l'interfaccia \texttt{ExplorationStrategy}; ognuna definisce
il metodo \texttt{next\_move} che restituisce la prossima cella verso cui spostarsi.
Tutte condividono una \textbf{priorit\`a universale}: se l'agente trasporta un oggetto,
la destinazione \`e il magazzino pi\`u vicino; se conosce un oggetto non reclamato, si dirige
verso di esso. Solo quando nessuna delle due condizioni \`e attiva la strategia di
esplorazione vera e propria viene invocata.

\section{Funzioni condivise della classe base}

La classe \texttt{ExplorationStrategy} espone sette metodi ereditati da tutte le strategie.

\paragraph{Priority-Move.}
Implementa la priorit\`a universale. Se l'agente trasporta un oggetto, lo indirizza
all'ingresso del magazzino scelto; altrimenti, se conosce oggetti non ancora reclamati,
lo dirige verso il pi\`u vicino (Manhattan). Solo se nessuna condizione \`e attiva il
controllo passa alla logica di esplorazione specifica.

\paragraph{Selezione ingresso di consegna.}
Sceglie dinamicamente l'ingresso di consegna combinando distanza e congestione.
Il costo per un ingresso $e$ \`e:
$$C(e) \;=\; d(\text{agent},\,e) \;+\; 4{,}5\cdot Q_{\text{local}}
         \;+\; 1{,}5\cdot R_{\text{reserved}} \;+\; \varepsilon_{\text{tie}}$$
dove $Q_{\text{local}}$ \`e il numero di agenti noti entro raggio~2 dall'ingresso,
$R_{\text{reserved}}$ \`e il numero di prenotazioni attive su quell'ingresso, e
$\varepsilon_{\text{tie}}$ \`e un tie-break deterministico basato
sull'ID dell'agente.
L'ingresso scelto viene \emph{prenotato} con un lock di 8~tick; un cambio
avviene solo se il nuovo costo \`e almeno il~20\% inferiore a quello corrente,
riducendo le oscillazioni.

\paragraph{Find-Frontier.}
Il metodo \texttt{ExplorationStrategy.\_find\_frontiers} individua le frontiere di
copertura: celle \texttt{EMPTY} gi\`a osservate dall'agente che confinano con almeno
una cella \texttt{EMPTY} non ancora vista nella mappa locale. Per ogni cella osservata
si verificano i quattro vicini cardinali; se almeno uno di essi \`e una cella
percorribile non presente in \texttt{seen\_cells}, la cella viene marcata come
frontiera. Il risultato viene memorizzato in cache per \texttt{agent\_id} e ricalcolato
solo quando aumenta il numero di celle in \texttt{seen\_cells}, riducendo il costo
computazionale nei cicli ripetuti di esplorazione.

\textit{Nota implementativa:} questo metodo era inizialmente stato sviluppato per testare
strategie con visione esclusivamente locale. Tuttavia, per mantenere omogeneita tra
tutti gli agenti e facilitare il confronto equo delle strategie, e stata successivamente
scartata a favore di un approccio unificato dove tutte le strategie accedono alla mappa
globale. Il metodo rimane nel codice come utility riutilizzabile per futuri sviluppi.

\paragraph{Coverage targets.}
Fornisce i target di esplorazione utilizzando la conoscenza globale della mappa
(accedendo direttamente a \texttt{env.grid.empty\_cells()}), con tre livelli di priorit\`a:
\begin{enumerate}
  \item Celle \texttt{EMPTY} mai osservate localmente (\textit{unexplored});
  \item Celle \texttt{EMPTY} osservate ma non visitate da almeno $\max(8,\,N/2)$ tick
        (\textit{stale});
  \item Tutte le celle \texttt{EMPTY} (\textit{patrolling}).
\end{enumerate}
Il passaggio al livello successivo avviene solo quando il precedente \`e vuoto. Questa
scelta di design consente alle strategie di operare con consapevolezza della struttura
complessiva del magazzino, migliorando l'efficienza esplorativa pur mantenendo autonomia
locale nelle decisioni di movimento.

\paragraph{Guadagno informativo (\textit{information gain}).}
Euristica condivisa da pi\'u strategie (Smart Random Walk e Greedy).
Per una cella candidata~$t$ il guadagno informativo \`e il numero di celle non ancora viste
che l'agente osserverebbe se si trovasse in~$t$:
$$\text{info\_gain}(t) = \left|\,\left\{(r,c) \in \mathcal{V}(t) : (r,c) \notin \text{seen\_cells}\right\}\,\right|$$
con
$$\mathcal{V}(t)=\left\{(r,c): |r-t_r|+|c-t_c|\le v_r\right\}.$$
Significato dei termini:
\begin{itemize}
  \item $t=(t_r,t_c)$: cella candidata verso cui l'agente sta valutando di muoversi;
  \item $r,c$: indici di riga e colonna di una cella della griglia;
  \item $|r-t_r|+|c-t_c|$: distanza di Manhattan tra la cella $(r,c)$ e il target $t$;
  \item $v_r$: raggio di visibilit\`a dell'agente; definisce il \emph{diamante} di celle
        osservabili da $t$ (insieme $\mathcal{V}(t)$);
  \item $\text{seen\_cells}$: insieme delle celle gi\`a osservate in precedenza
        dall'agente;
\end{itemize}
In pratica, $\text{info\_gain}(t)$ conta quanta nuova informazione produrrebbe il target
$t$: valori alti favoriscono zone con maggiore densit\`a di celle ancora inesplorate.

% ------------------------------------------------------------------
\section{Smart Random Walk}

Passeggiata casuale intelligente. Per ogni vicino percorribile viene calcolato un
punteggio che bilancia:

\begin{itemize}
  \item \textbf{Guadagno informativo} --- celle ignote nell'intorno del candidato,
        pesato da $w_{\text{info}} = 1{,}2$;
  \item \textbf{Stale bonus} --- tick dall'ultima visita alla cella, pesato da
        $w_{\text{stale}} = 0{,}8$;
  \item \textbf{Separazione} --- distanza media dagli agenti noti, pesata da
        $w_{\text{sep}} = 0{,}5$;
  \item \textbf{Penalit\`a rivisitazione} --- penalit\`a se la cella \`e stata visitata di recente,
        pesata da $w_{\text{revisit}} = 1{,}0$;
  \item \textbf{Rumore casuale} --- componente $\mathcal{U}(0,\, 0{,}3)$ per evitare pattern deterministici.
\end{itemize}

La mossa viene selezionata con probabilit\`a proporzionale al punteggio (\textit{softmax}),
mantenendo una componente stocastica che impedisce cicli.

\section{Frontier}

Esplorazione frontier-based classica. L'agente identifica le frontiere: celle percorribili
adiacenti ad almeno una cella \texttt{UNKNOWN} nella propria mappa locale. Tra tutte le
frontiere, seleziona quella pi\`u vicina per distanza di Manhattan e vi si dirige tramite A*.
Quando non esistono pi\`u frontiere (mappa completamente esplorata), l'agente esegue un
random walk tra i vicini liberi.

\section{Greedy}

Seleziona il target di esplorazione che minimizza un costo composito:

$$\text{cost}(t) = d(\text{agent}, t) + d(t, \text{wh}_{\text{nearest}}) - \text{info\_gain}(t)$$

dove $d$ \`e la distanza di Manhattan, $\text{wh}_{\text{nearest}}$ \`e il magazzino noto
pi\`u vicino e $\text{info\_gain}(t)$ conta le celle non ancora osservate entro il raggio
di visibilit\`a $v_r$ del target~$t$ (cfr.\ formula nella Sezione~3.1).
L'effetto \`e un bias verso zone inesplorate
e vicine ai magazzini. In assenza di target utili, l'agente si avvicina al peer pi\`u
vicino per favorire lo scambio di mappa.

\section{Sector}

La griglia viene suddivisa in bande orizzontali, una per agente (in base a
$\texttt{agent\_id} \bmod N$). Ogni agente copre sistematicamente la propria banda con
un pattern a serpentina, poi si sposta sulle bande adiacenti non coperte. Questo garantisce
copertura distribuita e uniforme senza bisogno di comunicazione esplicita.

\section{Repulsion}

Tra tutti i coverage target disponibili, seleziona quello che massimizza il punteggio:

$$\text{score}(t) = \frac{1}{|\mathcal{A}|}\sum_{a \in \mathcal{A}} d(t, a) - d(\text{agent}, t)$$

dove $\mathcal{A}$ \`e l'insieme delle posizioni degli agenti noti. La strategia massimizza
la distanza media dagli altri agenti (repulsione) e minimizza il costo di spostamento,
producendo una dispersione uniforme dello sciame sulle zone inesplorate.


